\documentclass[11pt]{article}
% Shared local declarations. Document roots use the repository's business-admin family.
\usepackage{business-admin}
\usepackage{amsmath,amssymb,tabularx,longtable,array,booktabs,enumitem}
\usepackage{ragged2e}
\newcolumntype{Y}{>{\RaggedRight\arraybackslash}X}
\newcommand{\SuiteMeta}[2]{\def\DocTitle{#1}\def\DocSubtitle{#2}\def\DocVersion{1.0}\def\DocStatus{Educational and operational reference}\def\DocOwner{Parliamentary Procedure and Governance Suite}\def\DocSummary{Independently authored reference and working materials, based on RONR 12th edition.}}
\newcommand{\SuiteOpen}{\makedoctitle\vspace{-0.5em}\noindent\rule{\linewidth}{0.4pt}\par\medskip}
\newcommand{\RonrAnchor}[1]{\par\smallskip{\footnotesize\textbf{Primary reading:} Robert's Rules of Order Newly Revised, 12th ed., #1. Original synthesis; not official interpretation.}\par}
\newcommand{\WriteLine}{\par\noindent\rule{\linewidth}{0.35pt}\par\vspace{0.5cm}}
\newcommand{\WriteLines}[1]{\par\noindent\rule{\linewidth}{0.35pt}\par\vspace{0.35cm}\noindent\rule{\linewidth}{0.35pt}\par\vspace{0.35cm}\noindent\rule{\linewidth}{0.35pt}\par\vspace{0.25cm}}
\newcommand{\Field}[1]{\par\medskip\noindent\textbf{#1}\par\WriteLines{3}}
\newcommand{\GovCheckbox}{\(\square\)}
\setlength{\emergencystretch}{3em}

\SuiteMeta{Operating Procedure: Control a Bylaws Amendment}{SOP | controlled meeting administration}
\begin{document}
\SuiteOpen
\RonrAnchor{56--57}
\section{Trigger and accountability}
\textbf{Trigger:} On introduction of governing-text change\quad\textbf{Responsible:} Secretary / bylaws committee / assembly

\textbf{Primary reading:} RONR (12th ed.) 56--57

\section{Operating steps}
\begin{enumerate}[leftmargin=*,itemsep=4pt]
\item Identify full current text and proposed amendment.
\item Apply notice and authorized-proposer rules.
\item Check any scope constraints on amendments.
\item Verify all voting and external-approval conditions.
\item Issue controlled adopted revision with effective-date record.
\end{enumerate}

\section{Quality gate}
The secretary or designated independent reviewer must be able to trace the action to a valid convening authority, applicable governing rule, eligible voters where relevant, the chosen disposition and an official record. A process pass does not establish statutory legality.

\section{Exception and escalation}
Proposal alters member rights beyond the notice scope. Document the issue, stop any unauthorized substantive action and obtain qualified parliamentary or legal guidance as appropriate.

\section{Retention and audit evidence}
\begin{itemize}[leftmargin=*,itemsep=4pt]
\item Approved meeting notice and governing authority where applicable.
\item Attendance and quorum evidence when a vote is involved.
\item Exact action text, vote basis and final disposition.
\item Corrections, chair rulings and approvals needed to reconstruct the transaction.
\end{itemize}
\end{document}
